\chapter{INTRODUCTION}

\section{Background}

Common bean (\textit{Phaseolus vulgaris} L.) is the most widely produced edible
legume crop, and dry bean seed is traded as a commodity whose price depends on
variety purity \cite{koklu2020}. Buyers distinguish Seker, Barbunya, Bombay,
Cali, Dermason, Horoz and Sira by size, shape and colour, and a consignment that
mixes varieties is downgraded. Grading is therefore performed at intake, usually
by trained inspectors who examine samples by hand. This is slow, subjective and
difficult to audit: two inspectors can disagree on beans that lie between two
variety descriptions.

Computer vision offers a deterministic alternative. A camera and a segmentation
step produce a binary mask for each grain, from which geometric descriptors such
as area, perimeter, axis lengths, eccentricity, solidity and roundness can be
computed in milliseconds. Köklü and Özkan built such a system for exactly the
seven varieties studied here and released both the measurement pipeline and the
resulting feature table as the UCI Dry Bean dataset \cite{uci602, koklu2020}.
Because the images have already been reduced to 16 numbers per bean, the
remaining task is a classical supervised learning problem rather than an image
recognition problem: no convolutional network is required, and a standard laptop
can reproduce every experiment in this report.

The learning problem is not trivial. Three of the seven varieties overlap
strongly in feature space, the classes are imbalanced by a factor of
\ImbalanceRatio, and several of the 16 descriptors are almost the same
measurement expressed twice. Deciding which classifier family handles those
conditions best, and demonstrating it with an evaluation protocol that does not
leak test information, is the substance of this case study.

\section{Problem Statement}

Develop a machine learning model to classify different varieties of dry beans
using shape, area, perimeter, roundness, compactness and other morphological
features extracted from bean images. Analyse feature distributions, perform
appropriate preprocessing and compare multiclass classification algorithms to
identify the most effective approach for distinguishing the seven dry bean
varieties.

The task is therefore a seven-class supervised classification problem with 16
continuous or integer inputs and one nominal target. The output required from the
model is a single variety label per bean, together with a class-probability
vector that quantifies the confidence of that label.

\section{Objectives}

\begin{enumerate}[label=\arabic*.]
  \item To study the morphological characteristics of the seven dry bean
        varieties and establish which descriptors carry between-class
        information.
  \item To validate the downloaded dataset and perform exploratory data analysis
        covering class distribution, feature distributions, correlation
        structure and outliers.
  \item To build a leakage-free preprocessing stage and a stratified
        training/test partition that is used unchanged by every experiment.
  \item To implement six multiclass classifiers spanning linear, instance-based,
        kernel and ensemble families, and to optimise their hyperparameters by
        stratified cross-validation on the training partition only.
  \item To compare the models on a common held-out test set using accuracy,
        macro-averaged precision, recall and F1, weighted F1, one-versus-rest
        ROC-AUC, confusion matrices, per-class scores and computational cost.
  \item To identify the best-performing algorithm from that evidence, explain
        why it succeeds, and analyse its errors and feature attributions.
  \item To demonstrate the selected model on representative beans through a
        reproducible prediction interface.
\end{enumerate}

\section{Scope of the Work}

The study covers the UCI Dry Bean dataset exactly as distributed: \DatasetRows{}
beans, 16 morphological features and seven varieties. It covers six classifier
families, systematic hyperparameter search, cross-validated model selection and a
single held-out evaluation, followed by error analysis and model interpretation.

The following are outside the scope. No raw images are processed; the work starts
from the published feature table, so segmentation quality and colour information
are not studied. No deep learning is attempted, because the inputs are 16
engineered scalars and a network would add cost without addressing the actual
difficulty, which is class overlap. No attempt is made to deploy the model in a
grading line, and no claim of industrial readiness is made: the evaluation is a
single random split of a single dataset collected in one laboratory. Additional
varieties, different cameras and different lighting are untested.

\section{Organization of the Report}

Chapter 2 reviews the original dry bean study and five further sources on kernel
selection, class imbalance, dimensionality reduction and deep alternatives, and
states the gap this case study addresses. Chapter 3 documents the dataset source,
licence and attributes and presents the exploratory analysis. Chapter 4 describes
data cleaning decisions, scaling, feature analysis and the split, including the
measures taken to prevent leakage. Chapter 5 develops the six algorithms with
their objective functions and hyperparameters and describes the implementation.
Chapter 6 reports the experimental results, compares the models, analyses the
errors and interprets the fitted model. Chapter 7 concludes, states the
limitations and lists realistic extensions. Appendix A documents the repository
and the reproduction steps; Appendix B carries the additional tables and figures.
